\documentclass[aspectratio=169,11pt]{beamer}
\usetheme{default}
\setbeamertemplate{navigation symbols}{}
\setbeamertemplate{footline}{%
  \hfill\insertframenumber/\inserttotalframenumber\hspace*{4pt}\vspace{2pt}%
}
\usepackage{lmodern}
\usepackage[most]{tcolorbox}
\usepackage{listings}
\usepackage{tikz}
\usetikzlibrary{arrows.meta,positioning,shapes.geometric,fit,backgrounds}
\usepackage{booktabs}
\usepackage{hyperref}

\definecolor{lcgreen}{HTML}{2E7D32}
\definecolor{lcblue}{HTML}{1565C0}
\definecolor{lcorange}{HTML}{EF6C00}
\definecolor{lcred}{HTML}{C62828}
\definecolor{codebg}{HTML}{F5F5F5}
\definecolor{docgray}{gray}{0.55}

\newtcolorbox{conceptbox}[1][]{colback=yellow!20,colframe=yellow!60!black,title=#1,fonttitle=\bfseries}
\newtcolorbox{warningbox}[1][]{colback=red!15,colframe=red!60!black,title=#1,fonttitle=\bfseries}
\newtcolorbox{tipbox}[1][]{colback=green!10,colframe=green!50!black,title=#1,fonttitle=\bfseries}

\newcommand{\docref}[1]{{\scriptsize\color{docgray}[Docs: #1]}}

\title[LangChain RAG]{Section 09: Retrieval-Augmented Generation (RAG)}
\author{LangChain Framework Guide}
\date{March 2026 --- LangChain v1.2.x}

\begin{document}

\begin{frame}
  \titlepage
\end{frame}

\begin{frame}{Agenda}
  \begin{itemize}
    \item RAG Mental Model: Giving LLMs access to your data
    \item RAG Fundamentals: Definition, motivation, pipeline
    \item Building RAG with LCEL: Complete working examples
    \item Advanced RAG Techniques: Multi-query, conversational, citations
    \item Evaluation and Production: Measuring quality, deployment checklist
    \item Best Practices: Common pitfalls, RAG vs fine-tuning
  \end{itemize}
  \vspace{1cm}
  \textcolor{lcblue}{Goal: Master practical RAG implementation from basics to production.}
\end{frame}

\begin{frame}{Mental Model: RAG = Access to Your Data}
  \begin{conceptbox}[Core Insight]
    RAG gives LLMs access to external knowledge without retraining.
    Think: ``LLM + Search + Context Injection''
  \end{conceptbox}
  \vspace{0.5cm}

  \begin{columns}[T]
    \column{0.48\textwidth}
    \textbf{Without RAG:}
    \begin{itemize}
      \item LLM knowledge frozen at training time
      \item Cannot answer questions about proprietary data
      \item Must fine-tune to add new knowledge
    \end{itemize}

    \column{0.48\textwidth}
    \textbf{With RAG:}
    \begin{itemize}
      \item Query retrieves relevant documents at runtime
      \item Context injected into prompt
      \item LLM answers based on retrieved docs
    \end{itemize}
  \end{columns}
\end{frame}

\begin{frame}{What is RAG? Definition and Motivation}
  \textbf{\textcolor{lcblue}{Retrieval-Augmented Generation (RAG):}}
  A technique combining information retrieval with generative AI to produce answers grounded in external documents.

  \begin{columns}[T]
    \column{0.5\textwidth}
    \textbf{Key Advantages:}
    \begin{itemize}
      \item No retraining needed
      \item Handles proprietary data
      \item Reduces hallucination
      \item Transparent sources
      \item Easy to update knowledge
    \end{itemize}

    \column{0.5\textwidth}
    \textbf{Why RAG Over Fine-tuning:}
    \begin{itemize}
      \item No GPU training required
      \item Knowledge stays current
      \item Can cite sources
      \item Works with large datasets
      \item Lower cost and faster deployment
    \end{itemize}
  \end{columns}
\end{frame}

\begin{frame}{The RAG Pipeline: Three Steps}
  \begin{itemize}
    \item[\textbf{1. Index}] Split documents into chunks, embed, store in vector DB
    \item[\textbf{2. Retrieve}] Find K most relevant chunks for query
    \item[\textbf{3. Generate}] Pass query with context to LLM, get answer
  \end{itemize}
  \vspace{0.8cm}

  \begin{center}
    Documents $\rightarrow$ Chunks $\rightarrow$ Embeddings $\rightarrow$ Vector DB
    \vspace{0.4cm}

    Query $\rightarrow$ Top-K Retrieval $\rightarrow$ LLM $\rightarrow$ Answer
  \end{center}

  \docref{docs.smith.langchain.com/rag}
\end{frame}

\begin{frame}[fragile]{Building a Basic RAG Chain (LCEL)}
  \fontsize{7}{8}\selectfont
  \texttt{
  from langchain.text\_splitter import RecursiveCharacterTextSplitter \\
  from langchain\_openai import OpenAIEmbeddings, ChatOpenAI \\
  from langchain\_community.vectorstores import Chroma \\
  from langchain.chains import create\_retrieval\_chain \\
  \\
  documents = load\_documents("./docs") \\
  splitter = RecursiveCharacterTextSplitter(chunk\_size=1000) \\
  chunks = splitter.split\_documents(documents) \\
  \\
  embeddings = OpenAIEmbeddings() \\
  vector\_db = Chroma.from\_documents(chunks, embeddings) \\
  retriever = vector\_db.as\_retriever(k=3) \\
  \\
  llm = ChatOpenAI(model="gpt-4") \\
  rag\_chain = create\_retrieval\_chain(retriever, doc\_chain) \\
  result = rag\_chain.invoke(\{"input": "What is RAG?"\}) \\
  }
  \normalsize
  \vspace{0.3cm}
  \docref{LangChain LCEL docs}
\end{frame}

\begin{frame}{The RAG Prompt: Formatting Context}
  \begin{tipbox}[Best Practice]
    Tell the LLM explicitly to use context and cite sources.
  \end{tipbox}

  \textbf{Good RAG Prompt Template:}
  \begin{itemize}
    \item Clear instruction to use context
    \item Role/expertise framing
    \item Clear separation: context, query, answer
    \item Fallback behavior for missing context
  \end{itemize}

  \textbf{Example Prompt:}
  ``You are an expert assistant. Answer based ONLY on the provided context. If context doesn't contain the answer, say 'Not in context.' Context: \{context\} Question: \{input\} Answer:''
\end{frame}

\begin{frame}{create\_stuff\_documents\_chain}
  \begin{conceptbox}[Stuff Strategy]
    Combine all retrieved docs into one prompt. Good for small/medium contexts.
  \end{conceptbox}

  \textbf{How it works:}
  \begin{enumerate}
    \item Retriever returns K documents
    \item All documents passed as single context block
    \item LLM generates response based on context
    \item Simple, effective, and fast
  \end{enumerate}

  \textbf{Advantages:}
  \begin{itemize}
    \item Minimal latency
    \item Good for most use cases
    \item Easy to implement
  \end{itemize}

  \textbf{Limitations:}
  \begin{itemize}
    \item Context window can fill quickly with large docs
    \item No ranking or compression of results
  \end{itemize}
\end{frame}

\begin{frame}[fragile]{create\_retrieval\_chain: Full RAG}
  \textbf{Combines retriever + document chain into single pipeline.}

  \fontsize{7}{8}\selectfont
  \texttt{
  from langchain.chains import create\_retrieval\_chain \\
  from langchain\_openai import OpenAIEmbeddings \\
  \\
  embeddings = OpenAIEmbeddings() \\
  db = FAISS.load\_local("vector\_db", embeddings) \\
  retriever = db.as\_retriever(k=4) \\
  llm = ChatOpenAI() \\
  \\
  rag\_chain = create\_retrieval\_chain(retriever, doc\_chain) \\
  \\
  result = rag\_chain.invoke(\{"input": "How does RAG work?"\}) \\
  print(result["answer"]) \\
  print(result["context"])
  }
  \normalsize
\end{frame}

\begin{frame}{Conversational RAG: Adding Chat History}
  \begin{warningbox}[Challenge]
    Standard RAG treats each query independently. Need to reformulate queries using conversation history.
  \end{warningbox}

  \textbf{Problem Example:}
  \begin{enumerate}
    \item User: ``Tell me about RAG''
    \item User: ``What about fine-tuning?''
    \item Without context, system retrieves about fine-tuning in general
    \item With context, system retrieves about fine-tuning vs RAG
  \end{enumerate}

  \begin{tipbox}[Solution]
    Use \texttt{create\_history\_aware\_retriever} to reformulate queries with chat history.
  \end{tipbox}
\end{frame}

\begin{frame}{Multi-Query RAG: Multiple Perspectives}
  \begin{conceptbox}[Intuition]
    One query formulation may miss relevant docs. Generate multiple query variations for better recall.
  \end{conceptbox}

  \begin{columns}[T]
    \column{0.5\textwidth}
    \textbf{Standard Query:}
    \begin{itemize}
      \item Single embedding
      \item May miss synonyms
      \item Limited perspective
    \end{itemize}

    \column{0.5\textwidth}
    \textbf{Multi-Query:}
    \begin{itemize}
      \item Generate 3-5 query variants
      \item Retrieve from each
      \item Combine and deduplicate
      \item Better coverage
    \end{itemize}
  \end{columns}

  \textbf{Example Variants for ``How do I build RAG?'':}
  \begin{itemize}
    \item ``Steps to implement retrieval-augmented generation''
    \item ``RAG architecture and setup guide''
    \item ``Creating a RAG system from scratch''
  \end{itemize}
\end{frame}

\begin{frame}{RAG with Sources: Returning Retrieved Docs}
  \begin{tipbox}[Transparency]
    Include source documents in response so users can verify answers.
  \end{tipbox}

  \textbf{Output Structure:}
  \begin{itemize}
    \item Answer: ``RAG is a technique that...''
    \item Context/Sources:
    \begin{itemize}
      \item docs/rag\_intro.pdf page 2
      \item docs/langchain\_guide.md section 5
      \item docs/tutorials.html
    \end{itemize}
  \end{itemize}

  \textbf{Benefits:}
  \begin{itemize}
    \item Users can verify answers
    \item Build trust in system
    \item Easy debugging
  \end{itemize}
\end{frame}

\begin{frame}{RAG with Citations: Making LLM Cite Passages}
  \begin{conceptbox}[Goal]
    Have LLM cite specific passages from retrieved documents.
  \end{conceptbox}

  \textbf{Citation Prompt Instruction:}
  \begin{itemize}
    \item ``Include [Citation: source\_name, passage] when citing''
    \item ``For each claim, provide the source document''
    \item ``Format as [Source: filename, section X]''
  \end{itemize}

  \textbf{Example Output:}
  \begin{itemize}
    \item ``RAG is [Citation: langchain\_docs, section 2.1] a technique that combines retrieval with generation.''
  \end{itemize}

  \textbf{Advantages:}
  \begin{itemize}
    \item Verifiable answers
    \item Transparent reasoning
    \item Reduced hallucination
  \end{itemize}
\end{frame}

\begin{frame}{Contextual Compression: Filter for Relevance}
  \begin{warningbox}[Problem]
    Retrieved chunks may contain irrelevant sections, wasting context window and confusing LLM.
  \end{warningbox}

  \begin{conceptbox}[Solution]
    Use LLM or embedding-based filter to keep only relevant passages.
  \end{conceptbox}

  \textbf{Approaches:}
  \begin{enumerate}
    \item LLM-based: Use LLM to score and filter documents
    \item Embedding-based: Compute similarity to query
    \item Keyword matching: Keep paragraphs with query terms
  \end{enumerate}

  \textbf{Benefits:}
  \begin{itemize}
    \item Fit more documents in context window
    \item Reduce LLM confusion
    \item Lower token usage and costs
  \end{itemize}
\end{frame}

\begin{frame}{Parent Document Retriever: Chunk vs. Context}
  \begin{conceptbox}[Idea]
    Index small chunks for fine-grained retrieval, but return larger parent docs for context.
  \end{conceptbox}

  \textbf{Problem Solved:}
  \begin{itemize}
    \item Small chunks: precise retrieval, no context
    \item Large chunks: good context, bad precision
    \item Solution: retrieve by chunk, return full section
  \end{itemize}

  \textbf{Configuration:}
  \begin{itemize}
    \item Child splitter: 100-500 tokens (search units)
    \item Parent splitter: 1500-2000 tokens (return units)
  \end{itemize}

  \textbf{Benefits:}
  \begin{itemize}
    \item Best of both worlds
    \item Better context for LLM
    \item Improved answer quality
  \end{itemize}
\end{frame}

\begin{frame}{RAG Evaluation: Measuring Quality}
  \begin{tipbox}[Metrics Matter]
    Evaluate RAG across multiple dimensions for production readiness.
  \end{tipbox}

  \begin{columns}[T]
    \column{0.48\textwidth}
    \textbf{Retrieval Quality:}
    \begin{itemize}
      \item Recall: Found all relevant docs?
      \item Precision: Were docs relevant?
      \item MRR: Rank of first match?
    \end{itemize}

    \column{0.48\textwidth}
    \textbf{Generation Quality:}
    \begin{itemize}
      \item BLEU/ROUGE: Text similarity
      \item Faithfulness: Grounded in context?
      \item Relevance: Matches query?
    \end{itemize}
  \end{columns}

  \begin{columns}[T]
    \column{0.48\textwidth}
    \textbf{Overall:}
    \begin{itemize}
      \item Hallucination rate
      \item Answer correctness
    \end{itemize}

    \column{0.48\textwidth}
    \textbf{Cost/Speed:}
    \begin{itemize}
      \item Latency (E2E)
      \item Tokens per query
      \item Dollar cost per query
    \end{itemize}
  \end{columns}
\end{frame}

\begin{frame}{RAG Fusion and Re-ranking Strategies}
  \begin{conceptbox}[Beyond Simple Retrieval]
    Combine multiple retrieval strategies and intelligently rank results.
  \end{conceptbox}

  \textbf{RAG Fusion:}
  \begin{itemize}
    \item Generate multiple queries
    \item Retrieve from each
    \item Reciprocal Rank Fusion (RRF) to combine
    \item Deduplicate and reorder
  \end{itemize}

  \textbf{Re-ranking Strategies:}
  \begin{itemize}
    \item LLM-based: Use LLM to score docs
    \item Cross-encoder: Specialized ranking model
    \item Diversity: Cover different aspects
    \item Semantic: Re-rank by relevance
  \end{itemize}

  \docref{langchain.retrievers}
\end{frame}

\begin{frame}{Streaming RAG Responses}
  \begin{tipbox}[User Experience]
    Stream LLM response token-by-token for faster perceived latency.
  \end{tipbox}

  \textbf{How Streaming Works:}
  \begin{enumerate}
    \item Retrieval happens upfront (not streamed)
    \item Only LLM generation is streamed
    \item Sources returned separately
    \item Improve perceived responsiveness
  \end{enumerate}

  \textbf{Considerations:}
  \begin{itemize}
    \item Latency still includes retrieval time
    \item Frontend must handle streaming
    \item Good for long responses
  \end{itemize}

  \begin{warningbox}[Limitation]
    Cannot stream documents; must retrieve first, then stream answer.
  \end{warningbox}
\end{frame}

\begin{frame}{Common RAG Pitfalls}
  \begin{warningbox}[Pitfall 1: Poor Chunking]
    Chunks too large lose specificity, too small lack context. Fix: 500-1500 tokens, 10-20 percent overlap.
  \end{warningbox}

  \begin{warningbox}[Pitfall 2: Stale Vector DB]
    Documents change but embeddings not updated. Fix: Re-index regularly, use delta updates.
  \end{warningbox}

  \begin{warningbox}[Pitfall 3: Hallucination Despite Context]
    LLM ignores context or makes up facts. Fix: Grounding prompts, lower temperature.
  \end{warningbox}

  \begin{warningbox}[Pitfall 4: Retrieval Failures]
    Query phrased differently from docs. Fix: Multi-query, hybrid search, query expansion.
  \end{warningbox}
\end{frame}

\begin{frame}{RAG vs Fine-tuning: Comparison}
  \begin{center}
    \small
    \begin{tabular}{|l|c|c|}
      \hline
      \textbf{Dimension} & \textbf{RAG} & \textbf{Fine-tune} \\
      \hline
      Training & No & Yes (GPU) \\
      Knowledge Update & Easy & Expensive \\
      Latency & Higher & Lower \\
      Citations & Easy & Hard \\
      Hallucination & Reduced & Possible \\
      Capacity & Large & Limited \\
      Deploy Cost & Lower & Higher \\
      Privacy & External DB & Model weights \\
      \hline
    \end{tabular}
  \end{center}

  \begin{tipbox}[When to Use RAG]
    Frequently updated knowledge, need citations, large datasets.
  \end{tipbox}

  \begin{tipbox}[When to Use Fine-tuning]
    Task-specific behavior, latency-critical, small datasets.
  \end{tipbox}
\end{frame}

\begin{frame}{Production RAG Checklist}
  \begin{conceptbox}[Before Launch]
    Ensure robust, measurable, and maintainable system.
  \end{conceptbox}

  \begin{columns}[T]
    \column{0.5\textwidth}
    \textbf{Data and Indexing:}
    \begin{itemize}
      \item Documents loaded and versioned
      \item Chunking strategy tested
      \item Vector DB production-ready
      \item Embeddings model chosen
    \end{itemize}

    \column{0.5\textwidth}
    \textbf{Evaluation and Testing:}
    \begin{itemize}
      \item Human eval on sample set
      \item Automated metrics (RAGAS)
      \item A/B testing framework
      \item Baseline established
    \end{itemize}
  \end{columns}

  \begin{columns}[T]
    \column{0.5\textwidth}
    \textbf{Operations:}
    \begin{itemize}
      \item Caching strategy
      \item Rate limiting
      \item Error handling
      \item Logging and audit trail
    \end{itemize}

    \column{0.5\textwidth}
    \textbf{Monitoring:}
    \begin{itemize}
      \item Latency tracked
      \item Token costs logged
      \item Hallucination detection
      \item User feedback loop
    \end{itemize}
  \end{columns}
\end{frame}

\begin{frame}{Advanced RAG: Adaptive Routing}
  \begin{conceptbox}[Routing Based on Query]
    Use router to decide: retrieve or answer directly? Which retriever?
  \end{conceptbox}

  \textbf{Routing Strategies:}
  \begin{itemize}
    \item Query classification: Does it need RAG?
    \item Retriever selection: Keyword vs semantic vs hybrid
    \item Document source routing: Which DB?
    \item Confidence-based: Use fallback if uncertain
  \end{itemize}

  \textbf{Benefits:}
  \begin{itemize}
    \item Reduced latency for simple queries
    \item Better results for complex queries
    \item Cost optimization
  \end{itemize}
\end{frame}

\begin{frame}{Caching in RAG: Speed and Cost}
  \begin{tipbox}[Opportunities]
    Cache at multiple levels: embeddings, retrievals, responses.
  \end{tipbox}

  \begin{columns}[T]
    \column{0.5\textwidth}
    \textbf{Query Caching:}
    \begin{itemize}
      \item Cache exact same queries
      \item FIFO or LRU cache
      \item Limited by diversity
    \end{itemize}

    \column{0.5\textwidth}
    \textbf{Semantic Caching:}
    \begin{itemize}
      \item Cache by query intent
      \item Return similar answers
      \item More flexible
    \end{itemize}
  \end{columns}

  \textbf{Implementation:}
  \begin{itemize}
    \item In-memory cache for development
    \item Redis for production
    \item LangChain built-in support
  \end{itemize}
\end{frame}

\begin{frame}{Hybrid Search: Keyword plus Semantic}
  \begin{conceptbox}[Best of Both Worlds]
    Combine exact keyword matching with semantic similarity.
  \end{conceptbox}

  \begin{columns}[T]
    \column{0.48\textwidth}
    \textbf{Keyword Search:}
    \begin{itemize}
      \item BM25, TF-IDF, Lucene
      \item Exact term matching
      \item Good for names/acronyms
      \item Fast, no embeddings
    \end{itemize}

    \column{0.48\textwidth}
    \textbf{Semantic Search:}
    \begin{itemize}
      \item Vector embeddings
      \item Synonym understanding
      \item Good for concepts
      \item Slower, requires embeddings
    \end{itemize}
  \end{columns}

  \textbf{Approach:}
  \begin{itemize}
    \item Retrieve from both
    \item Rank and deduplicate
    \item Return combined results
  \end{itemize}
\end{frame}

\begin{frame}{Metadata Filtering in RAG}
  \begin{tipbox}[Pre-filtering]
    Use metadata to restrict retrieval before semantic search.
  \end{tipbox}

  \textbf{Common Metadata:}
  \begin{itemize}
    \item Date: Only 2025 onward
    \item Source: Official docs only
    \item Category: API reference only
    \item Language: English only
    \item Author: Peer-reviewed sources
  \end{itemize}

  \textbf{Benefits:}
  \begin{itemize}
    \item Filter irrelevant results early
    \item Reduce context noise
    \item Improve answer quality
    \item Lower costs
  \end{itemize}
\end{frame}

\begin{frame}{Chain of Thought RAG}
  \begin{conceptbox}[Idea]
    Combine multi-step reasoning with iterative retrieval.
  \end{conceptbox}

  \textbf{Pattern:}
  \begin{enumerate}
    \item LLM breaks complex question into sub-questions
    \item Retrieve context for each sub-question
    \item Combine answers into final response
  \end{enumerate}

  \textbf{Example:}
  \begin{itemize}
    \item Q: ``Compare RAG and fine-tuning for proprietary documents''
    \item Sub-Q1: ``What is RAG?''
    \item Sub-Q2: ``What is fine-tuning?''
    \item Sub-Q3: ``Advantages/disadvantages of each?''
  \end{itemize}

  \textbf{Benefit:} Better answers to complex, multi-faceted questions
\end{frame}

\begin{frame}{Dynamic Context Window Management}
  \begin{warningbox}[Challenge]
    Limited context window. Too much: LLM confused. Too little: quality suffers.
  \end{warningbox}

  \textbf{Strategies:}
  \begin{itemize}
    \item Adaptive K: Increase K if low confidence
    \item Token Budget: Reserve tokens for response
    \item Truncation: Summarize long docs
    \item Hierarchical: Summaries then full docs
  \end{itemize}

  \textbf{Implementation:}
  \begin{itemize}
    \item Monitor token usage
    \item Adjust K based on quality
    \item Track latency impact
  \end{itemize}
\end{frame}

\begin{frame}{Testing and Debugging RAG}
  \textbf{Testing Strategy:}
  \begin{itemize}
    \item Unit tests: Chunking, embeddings, indexing
    \item Integration tests: Retrieval accuracy, prompt formatting
    \item E2E tests: Full pipeline, known Q\&A pairs
    \item Regression tests: Catch improvements and regressions
  \end{itemize}

  \textbf{Debugging Checklist:}
  \begin{itemize}
    \item Poor answers: Check retrieved docs, prompt, LLM
    \item Slow queries: Profile retrieval vs LLM time
    \item Out of memory: Reduce docs, chunk size
    \item Hallucinations: Improve prompts, lower temperature
  \end{itemize}

  \textbf{Tools:}
  \begin{itemize}
    \item LangSmith for tracing
    \item RAGAS for metrics
    \item Custom evaluation datasets
  \end{itemize}
\end{frame}

\begin{frame}{RAG Security: PII and Sensitive Data}
  \begin{warningbox}[Risk]
    RAG can expose PII if documents contain sensitive information.
  \end{warningbox}

  \textbf{Mitigation Strategies:}
  \begin{itemize}
    \item PII Redaction: Remove/mask emails, SSNs before indexing
    \item Access Control: Only authorized users see certain docs
    \item Query Filtering: Restrict sensitive field retrieval
    \item Output Filtering: Post-process LLM output
  \end{itemize}

  \textbf{Best Practices:}
  \begin{itemize}
    \item Audit documents before indexing
    \item Use metadata filtering
    \item Monitor for data leakage
    \item Test with sensitive queries
  \end{itemize}
\end{frame}

\begin{frame}{RAG Cost Optimization}
  \begin{tipbox}[Budget-Conscious Approaches]
    Reduce embedding, LLM, and DB costs.
  \end{tipbox}

  \textbf{Optimization Strategies:}
  \begin{enumerate}
    \item Smaller Models: Cheaper embeddings, smaller LLMs
    \item Batch Processing: Index once, reuse embeddings
    \item Reduce K: Start with 3 instead of 10
    \item Caching: Cache embeddings and responses
    \item Compression: Summarize docs to fit more
    \item Approximate Search: HNSW, IVF vs exact
  \end{enumerate}

  \textbf{Typical Costs:}
  \begin{itemize}
    \item RAG query: dollars 0.001--dollars 0.01
    \item Depends on model, context size, vector DB
  \end{itemize}
\end{frame}

\begin{frame}{Knowledge Graph plus RAG}
  \begin{conceptbox}[Hybrid Approach]
    Combine semantic search with structured knowledge graphs.
  \end{conceptbox}

  \textbf{Pattern:}
  \begin{itemize}
    \item Extract entities and relationships from docs
    \item Build knowledge graph (Neo4j, RDF)
    \item Query graph for structured context
    \item Pass graph plus retrieved docs to LLM
  \end{itemize}

  \textbf{Benefits:}
  \begin{itemize}
    \item Better reasoning on relationships
    \item More precise, entity-based retrieval
    \item Combines semantic plus structured search
  \end{itemize}
\end{frame}

\begin{frame}{Multi-Modal RAG: Images plus Text}
  \begin{conceptbox}[Beyond Text]
    Extend RAG to images, videos, audio.
  \end{conceptbox}

  \textbf{Challenges:}
  \begin{itemize}
    \item No standard multi-modal embeddings yet
    \item Slower inference, larger models
    \item Complex indexing (text plus images)
  \end{itemize}

  \textbf{Approaches:}
  \begin{itemize}
    \item Vision-Language Models: CLIP, LLaVA
    \item Hybrid Indexing: Store image embeddings and descriptions
    \item Image-to-Text: Captions then standard RAG
  \end{itemize}

  \textit{Emerging area with use case-dependent solutions.}
\end{frame}

\begin{frame}{RAG Cheat Sheet: Quick Reference}
  \textbf{Basic RAG in 3 lines:}
  \begin{itemize}
    \item \texttt{retriever = db.as\_retriever(k=3)}
    \item \texttt{rag\_chain = create\_retrieval\_chain(retriever, doc\_chain)}
    \item \texttt{result = rag\_chain.invoke(\{``input'': question\})}
  \end{itemize}

  \textbf{Key Components:}
  \begin{itemize}
    \item Retriever: vector DB, keyword, multi-query
    \item Doc Chain: stuff, map-reduce, refine
    \item Full Chain: retrieval, history-aware
  \end{itemize}

  \textbf{Production Considerations:}
  \begin{itemize}
    \item Evaluate with RAGAS
    \item Monitor with LangSmith
    \item Cache aggressively
    \item Test thoroughly
  \end{itemize}
\end{frame}

\begin{frame}{Self-Check Questions}
  \begin{enumerate}
    \item What is RAG and how does it differ from fine-tuning?
    \item Explain the three steps of a RAG pipeline.
    \item How does multi-query retrieval improve recall?
    \item What is contextual compression and why use it?
    \item How do you add chat history to RAG?
    \item Name three RAG evaluation metrics.
    \item What are common RAG pitfalls and fixes?
    \item How can you reduce RAG costs?
  \end{enumerate}

  \vspace{0.5cm}
  \textcolor{lcorange}{\textbf{Hint:}} Review slides on basics (frames 4-5), LCEL (frames 6-9), and evaluation (frame 18).
\end{frame}

\begin{frame}{Summary and Next Steps}
  \begin{conceptbox}[Key Takeaway]
    RAG is the practical way to give LLMs access to your data without retraining.
  \end{conceptbox}

  \textbf{What You've Learned:}
  \begin{itemize}
    \item RAG fundamentals and mental model
    \item Building RAG chains with LCEL
    \item Advanced techniques: multi-query, compression, citations
    \item Evaluation and production best practices
    \item Common pitfalls and optimization strategies
  \end{itemize}

  \textbf{Next Steps:}
  \begin{enumerate}
    \item Build a simple RAG pipeline with your documents
    \item Evaluate using RAGAS or custom metrics
    \item Iterate on retrieval (chunk size, K, retriever)
    \item Deploy with monitoring and caching
    \item Collect feedback and improve continuously
  \end{enumerate}
\end{frame}

\begin{frame}{Resources and Further Reading}
  \textbf{Official LangChain Docs:}
  \begin{itemize}
    \item \url{https://docs.smith.langchain.com/rag}
    \item \url{https://docs.smith.langchain.com/rag-from-scratch}
    \item \url{https://api.python.langchain.com}
  \end{itemize}

  \textbf{Evaluation and Tracing:}
  \begin{itemize}
    \item RAGAS: \url{https://github.com/explodinggradients/ragas}
    \item LangSmith: \url{https://docs.smith.langchain.com}
  \end{itemize}

  \textbf{Research:}
  \begin{itemize}
    \item Retrieval-Augmented Generation for Knowledge-Intensive NLP Tasks (Lewis et al., 2020)
    \item In-Context Retrieval-Augmented Language Models (Shi et al., 2023)
  \end{itemize}
\end{frame}

\end{document}
